\documentclass[12pt,a4paper]{article}
\usepackage[spanish]{babel}
\usepackage[utf8]{inputenc}
\usepackage[T1]{fontenc}
\usepackage{lmodern}
\usepackage[table]{xcolor}
\usepackage{geometry}
\usepackage{setspace}
\usepackage{longtable}
\usepackage{booktabs}
\usepackage{array}
\usepackage{tabularx}
\usepackage{hyperref}
\usepackage{enumitem}

\geometry{margin=2.4cm}
\setstretch{1.12}
\setlength{\parindent}{0pt}
\setlength{\parskip}{0.55em}
\hypersetup{colorlinks=true, linkcolor=black, urlcolor=blue, citecolor=black}

\definecolor{PrimaryBlue}{HTML}{16324F}
\definecolor{SoftGray}{HTML}{F3F6F8}

	itle{\Huge\textbf{Especificaci\'on de Requisitos de Software}\[0.3em]\Large Plataforma Web de Mascotas 3D}
\author{Bryan Quispe}
\date{10 de mayo de 2026}

\begin{document}

\begin{titlepage}
\centering
\vspace*{1.8cm}
{\Huge\bfseries\color{PrimaryBlue} ERS\\[0.35em]}
{\LARGE Especificaci\'on de Requisitos de Software\\[0.5em]}
{\large Plataforma Web de Mascotas 3D\\[1.2cm]}

\begin{tabular}{p{5cm} p{8cm}}
\rowcolor{SoftGray}
\textbf{Versi\'on} & 1.0 \\
\textbf{Fecha} & 10 de mayo de 2026 \\
\textbf{Autor} & Bryan Quispe \\
\textbf{Estado} & En desarrollo \\
\textbf{Alcance} & MVP en desarrollo: Autenticaci\'on, mascotas y visualizaci\'on 3D base \\
\end{tabular}

\vfill
{\large Documento de requisitos para compilaci\'on en Overleaf}
\end{titlepage}

\tableofcontents
\newpage

\section{Introducci\'on}
Este documento especifica los requisitos funcionales y no funcionales de la plataforma web de mascotas 3D en desarrollo. Ser\'a la base para el dise\~no, desarrollo y validaci\'on de la soluci\'on.

\section{Requisitos Funcionales}

\subsection{RF-01: Registro de usuario}
\begin{itemize}[leftmargin=*]
    \item \textbf{Descripci\'on}: La plataforma debe permitir el registro de nuevos usuarios.
    \item \textbf{Criterios de aceptaci\'on}:
    \begin{itemize}[leftmargin=*]
        \item Validar email \'unico en la base de datos.
        \item Contrase\~na debe tener al menos 8 caracteres.
        \item Contrase\~na debe estar hasheada con Bcrypt (10 rondas).
        \item Mensaje de error claro si email existe.
        \item Respuesta REST 201 Created en caso de \'exito.
    \end{itemize}
    \item \textbf{Prioridad}: Alta
    \item \textbf{Sprint}: 1
\end{itemize}

\subsection{RF-02: Inicio de sesi\'on (Login)}
\begin{itemize}[leftmargin=*]
    \item \textbf{Descripci\'on}: La plataforma debe permitir que usuarios registrados inicien sesi\'on y accedan a rutas protegidas.
    \item \textbf{Criterios de aceptaci\'on}:
    \begin{itemize}[leftmargin=*]
        \item Validar email y contrase\~na contra la base de datos.
        \item Generar JWT con identificador de usuario y fecha de expiraci\'on.
        \item Respuesta REST 200 OK con token en el body.
        \item Mensaje gen\'erico si las credenciales no coinciden.
    \end{itemize}
    \item \textbf{Prioridad}: Alta
    \item \textbf{Sprint}: 1
\end{itemize}

\subsection{RF-03: Registro de mascota}
\begin{itemize}[leftmargin=*]
    \item \textbf{Descripci\'on}: Usuario autenticado puede crear un registro de mascota.
    \item \textbf{Criterios de aceptaci\'on}:
    \begin{itemize}[leftmargin=*]
        \item Campos obligatorios: nombre, edad, tipo de animal.
        \item Asociar mascota al usuario propietario.
        \item Timestamp de creaci\'on autom\'atico.
        \item Respuesta REST 201 Created.
    \end{itemize}
    \item \textbf{Prioridad}: Alta
    \item \textbf{Sprint}: 1
\end{itemize}

\subsection{RF-04: Listado de mascotas}
\begin{itemize}[leftmargin=*]
    \item \textbf{Descripci\'on}: Usuario autenticado puede ver la lista de sus mascotas.
    \item \textbf{Criterios de aceptaci\'on}:
    \begin{itemize}[leftmargin=*]
        \item Retornar solo mascotas asociadas al usuario autenticado.
        \item Incluir nombre, edad, tipo y fecha de creaci\'on.
        \item Respuesta REST 200 OK con lista JSON.
    \end{itemize}
    \item \textbf{Prioridad}: Alta
    \item \textbf{Sprint}: 1
\end{itemize}

\subsection{RF-05: Visualizaci\'on de detalle de mascota y modelo 3D}
\begin{itemize}[leftmargin=*]
    \item \textbf{Descripci\'on}: Usuario autenticado puede ver el detalle de una mascota y su modelo 3D base.
    \item \textbf{Criterios de aceptaci\'on}:
    \begin{itemize}[leftmargin=*]
        \item Mostrar datos de la mascota y foto asociada si existe.
        \item Cargar modelo 3D base seg\'un el tipo de animal.
        \item Permitir rotaci\'on b\'asica del modelo en el visor.
    \end{itemize}
    \item \textbf{Prioridad}: Alta
    \item \textbf{Sprint}: 2
\end{itemize}

\subsection{RF-06: Carga de fotos de mascota}
\begin{itemize}[leftmargin=*]
    \item \textbf{Descripci\'on}: Usuario puede subir foto de su mascota.
    \item \textbf{Criterios de aceptaci\'on}:
    \begin{itemize}[leftmargin=*]
        \item Aceptar JPG y PNG.
        \item Tama\~no m\'aximo 5MB.
        \item Compresi\'on autom\'atica.
        \item Almacenamiento seguro en servidor.
    \end{itemize}
    \item \textbf{Prioridad}: Media
    \item \textbf{Sprint}: 2
\end{itemize}

\subsection{RF-07: Gestión de modelos 3D (Admin)}
\begin{itemize}[leftmargin=*]
    \item \textbf{Descripci\'on}: Un usuario con rol `admin` puede subir, versionar y administrar modelos 3D (GLB/GLTF). Adicionalmente, un usuario autenticado puede editar o personalizar versiones permitidas del modelo seg\'un las reglas del sistema.
    \item \textbf{Criterios de aceptaci\'on}:
    \begin{itemize}[leftmargin=*]
        \item El admin puede subir archivos GLB/GLTF y asignar metadatos (tipo de animal, nombre, versión).
        \item Versionado simple: cada subida puede crear una nueva versión del modelo.
        \item Permisos: solo admins pueden eliminar o publicar modelos para uso en el catálogo.
        \item El usuario autenticado solo puede editar campos permitidos o solicitar revisión cuando el cambio requiere aprobación.
        \item Respuesta REST 201 Created en subida exitosa.
    \end{itemize}
    \item \textbf{Prioridad}: Alta
    \item \textbf{Sprint}: 2
\end{itemize}

\section{Requisitos No Funcionales}

\subsection{RNF-01: Seguridad}
\begin{itemize}[leftmargin=*]
    \item Todas las contrase\~nas hasheadas con Bcrypt.
    \item JWT firmado con clave privada.
    \item CORS configurado para solo dominio propio.
    \item Sin SQL injection en queries.
\end{itemize}

\subsection{RNF-02: Rendimiento}
\begin{itemize}[leftmargin=*]
    \item Endpoint de login: <500ms.
    \item Listado de mascotas: <1000ms.
    \item Carga de modelo 3D: <2 segundos.
\end{itemize}

\subsection{RNF-03: Disponibilidad}
\begin{itemize}[leftmargin=*]
    \item Disponibilidad esperada: 95\% en MVP.
    \item Tiempo de recuperaci\'on ante fallo: <30 minutos.
\end{itemize}

\subsection{RNF-04: Compatibilidad}
\begin{itemize}[leftmargin=*]
    \item Backend: Node.js + NestJS.
    \item Frontend: Next.js + React.
    \item Base de datos: PostgreSQL.
    \item Navegadores: Chrome, Firefox, Edge \'ultimas versiones.
\end{itemize}

\subsection{RNF-05: Gestión de sesión / Timeout}
\begin{itemize}[leftmargin=*]
    \item La sesión del usuario expirará automáticamente tras 10 minutos de inactividad.
    \item Al cerrar la aplicación el cliente debe eliminar credenciales locales (token) para proteger la sesión.
    \item El servidor puede aplicar políticas de expiraci\'on o revocaci\'on de tokens para cumplir el requisito.
\end{itemize}

\section{Restricciones}
\begin{itemize}[leftmargin=*]
    \item Presupuesto limitado para hosting en MVP.
    \item Equipo de 3-4 desarrolladores.
    \item Plazo: 4 semanas para el MVP.
    \item Sin dependencias externas para modelos 3D base.
\end{itemize}

\section{Glosario}
\begin{itemize}[leftmargin=*]
    \item \textbf{JWT}: JSON Web Token para autenticaci\'on.
    \item \textbf{PBI}: Product Backlog Item.
    \item \textbf{MVP}: M\'inimo Producto Viable.
    \item \textbf{GLB}: Formato binario para modelos 3D.
\end{itemize}
\newpage
\section{Referencias}
Las siguientes fuentes se usaron como referencia para el formato y las convenciones de este documento:
\begin{thebibliography}{9}
\bibitem{scrumguide}
K. Schwaber y J. Sutherland,
\textit{The Scrum Guide: The Definitive Guide to Scrum: The Rules of the Game},
The Scrum Guide, 2020. Dispon\'ible en: \url{https://scrumguides.org}

\bibitem{ieee830}
IEEE,
\textit{IEEE Recommended Practice for Software Requirements Specifications},
IEEE Std 830-1998 (recomendaci\'on ampliamente utilizada para estructura de SRS).
\end{thebibliography}

\end{document}
